\section{Síntesis y ampliación}

\begin{enumerate}
    \item La clonación de comportamiento es la vía más rápida para comenzar, pero hay que estar atento a la conducta promedio y al covariate shift.
    \item Los modelos generativos dotan a la política de capacidad de expresión multimodal, y es necesario ajustar KL/temperature en la etapa de muestreo.
    \item DAgger/HG-DAgger/las estrategias híbridas de RL son una vía eficaz para resolver los compounding errors.
    \item El audit, los metadata y la colaboración multi-role de los datos de demostración de alta calidad son la base para la puesta en producción.
    \item La etapa de despliegue debe girar en torno al evidence board, vinculando las tres líneas de benchmark, safety y ops a la iteración de monitoreo.
    \item Hay que registrar cada evento de mitigation y rollback, para que el governance audit tenga fundamento.
\end{enumerate}

\subsection{Hoja de ruta de puesta en producción}

\begin{figure}[H]
\centering
\includegraphics[width=0.8\textwidth]{slide-02-14.jpg}
\caption{La diapositiva 02-14 muestra en forma de timeline los deliverables desde los datos hasta el deployment.}
\end{figure}

La hoja de ruta de puesta en producción divide el pipeline de aprendizaje por imitación en tres etapas: 1) preparación de datos (demo catalog + metadata), 2) entrenamiento de la política (diffusion/MoG + flow calibration), 3) recuperación ante desastres (drift playbook + governance review). Cada etapa debe definir guardrail, evaluation metric y artifacts.

\begin{table}[H]
\centering
\begin{tabular}{p{0.25\textwidth} p{0.55\textwidth}}
\toprule
Etapa & Key artifact \\
\midrule
Preparación de datos & Metadata-rich demo catalog + coverage report \\
Entrenamiento de la política & KL drift sweep + sampling temperature record \\
Recuperación de despliegue & Drift playbook + rollback checklist \\
\bottomrule
\end{tabular}
\caption{Hoja de ruta por etapas para la puesta en producción del aprendizaje por imitación}
\end{table}

\begin{warningbox}{No confundas el roadmap con un plan de proyecto}
El roadmap debe actualizarse constantemente: se obtienen el delay y el drift reales del evidence dashboard, y luego los lesson learned se escriben de vuelta en el timeline, para poder seguir avanzando.
\end{warningbox}

\subsection{Tabla de resumen}

\begin{table}[H]
\centering
\begin{tabular}{p{0.2\textwidth} p{0.35\textwidth} p{0.35\textwidth}}
\toprule
Tema & Takeaway central & Acción de ingeniería \\
\midrule
Clonación de comportamiento & Puesta en marcha rápida, pero hay que atender el covariance shift & Ejecutar un KL sweep después del entrenamiento y un sanity check antes del despliegue \\
Expresión de la política & La distribución generativa necesita explorarse por muestreo & Usar diffusion/MoG y registrar el sampling temperature \\
Corrección en línea & DAgger compensa los estados no vistos & Establecer un runbook de Observe-Assess-Act-Document \\
Control de calidad de datos & Multi-role audit & Usar metadata + clustering para identificar coverage gaps \\
Gobernanza del despliegue & Evidence dashboard & Vincular las tres líneas de benchmark/safety/ops \\
\bottomrule
\end{tabular}
\caption{Síntesis accionable de esta clase}
\end{table}

\subsection{Lecturas adicionales}
\begin{itemize}
    \item Chi et al., \textit{Diffusion Policy}: incorpora el diffusion model a las políticas robóticas.
    \item Ross et al., \textit{A Reduction of Imitation Learning and Structured Prediction to No-Regret Online Learning}: el artículo original de DAgger.
    \item Brohan et al., \textit{RT-2}: muestra los efectos del aprendizaje por imitación multimodal.
    \item Rajeswaran et al., \textit{Behavioral Cloning with Expert Interventions}: sistematiza las human interventions.
\end{itemize}

\subsection{Resumen de la sección}

Esta clase recorre desde la clonación de comportamiento hasta las políticas generativas, y desde el DAgger/online pipeline hasta el data QA y la gobernanza del despliegue, conformando un marco práctico end-to-end para el aprendizaje por imitación. Solo cuando todas las etapas están conectadas entre sí, el aprendizaje por imitación puede ejecutarse de manera verdaderamente confiable en el entorno de producción.

\end{document}
